\section*{Metodología}
\justifying

El planteamiento de las actividades se divide en las distintas etapas para cumplir
con los objetivos del proyecto:

\begin{enumerate}[noitemsep]
    \item \textbf{Selección de módulos:} Se seleccionarán los módulos a implementar
    en el sistema. Se tendrá en cuenta la importancia de los módulos en la enseñanza de la electrónica
    de potencia y su relevancia en la formación de los estudiantes. Se analizará el potencial
    didáctico e interactivo de cada módulo y su relación con los conceptos teóricos de la materia.

    \item \textbf{Diseño de prototipos:} Se diseñarán circuitos básicos de los módulos a implementar,
    con el fin de realizar pequeños prototipos y pruebas de concepto. Se analizarán los resultados
    de las pruebas y se realizarán ajustes en los circuitos, de ser necesario.

    \item \textbf{Diseño de módulos:} Una vez validados los prototipos, se procederá al diseño de los
    circuitos y placas definitivas de los módulos, así como un protocolo de comunicación entre el
    sistema y los módulos, para garantizar su correcto funcionamiento. Se seleccionarán los
    componentes electrónicos a utilizar, según cálculos de potencia y disipación de calor.
    Se diseñarán para que sean fáciles de intercambiar y conectar.
    
    \item \textbf{Diseño del sistema:} Se diseñará una interfaz de usuario intuitiva y accesible
    que permita a los usuarios interactuar con el sistema de forma sencilla. Se analizarán las dimensiones
    y la disposición de los módulos para garantizar la comodidad y la eficiencia en su uso. Se garantizará
    la seguridad del sistema y de los usuarios, mediante la implementación de medidas de protección.
    Además, se tendrá en cuenta el fácil mantenimiento y reparación del sistema en caso de necesitarlo.
    Se programará un microcontrolador para la interfaz de usuario y la comunicación con los módulos.
    
    \item \textbf{Construcción de los módulos y del sistema:} Una vez diseñados los módulos y el sistema,
    se procederá a la construcción de los mismos. Se realizarán pruebas de funcionamiento
    de los módulos y se analizarán los resultados. Se realizarán pruebas de funcionamiento del sistema
    en su totalidad, para garantizar su correcto funcionamiento.
    
    \item \textbf{Pruebas de funcionamiento:} Se realizarán pruebas del sistema en su totalidad para
    garantizar su correcto funcionamiento.
    
    \item \textbf{Documentación:} Se desarrollará un manual de usuario claro y detallado que explique
    el funcionamiento teórico de cada módulo. Se elaborará una guía práctica que oriente sobre los
    parámetros a modificar y los efectos a analizar.

    \item \textbf{Evaluación:} Se realizarán pruebas de funcionamiento del sistema con usuarios finales,
    con el fin de evaluar su desempeño y recopilar realimentación para futuras mejoras.
\end{enumerate}
